\documentclass[11pt]{article}

\usepackage{amsmath, amssymb}
\usepackage{geometry}
\geometry{margin=1in}

\title{\textbf{Lecture 8 Scribe}\\Gaussian, Uniform, Exponential, Laplace, and Gamma Random Variables}
\date{}
\author{}

\begin{document}
\maketitle

\section*{8. Gaussian, Uniform, Exponential, Laplace, and Gamma Random Variables}

\subsection*{8.1 Gaussian Random Variable — Definition}

\textbf{Definition.}
A Gaussian random variable is one whose PDF can be written in the general form
\[
f_X(x)=\frac{1}{\sqrt{2\pi\sigma^2}}
\exp\!\left(-\frac{(x-m)^2}{2\sigma^2}\right).
\]

\textbf{Notation.}
\[
X\sim \mathcal N(m,\sigma^2),
\]
where $m$ is the mean and $\sigma^2$ is the variance.

\subsection*{8.2 Expectation and Moments}

\textbf{Expected value of a function of $X$:}

\begin{itemize}
\item Continuous RV:
\[
E[g(X)]=\int g(x)f_X(x)\,dx.
\]

\item Discrete RV:
\[
E[g(X)]=\sum_k g(x_k)P_X(x_k).
\]
\end{itemize}

\textbf{Linearity:}
\[
E[aX+b]=aE[X]+b.
\]

For sums,
\[
E\!\left[\sum_{k=1}^{N} g_k(X)\right]
=\sum_{k=1}^{N}E[g_k(X)].
\]

\subsubsection*{Central Moments}

Second central moment (variance):
\[
E[(X-\mu_X)^2]=\sigma_X^2.
\]

Third central moment (skewness):
\[
C_s=\frac{E[(X-\mu_X)^3]}{\sigma_X^3}.
\]

Interpretation: measure of symmetry of the PDF about $\mu_X$.
Positive value $\rightarrow$ right.
Negative value $\rightarrow$ left.

Fourth central moment (kurtosis):
\[
C_k=\frac{E[(X-\mu_X)^4]}{\sigma_X^4}.
\]

Large value indicates a large peak near the mean.

\subsection*{8.3 Standard Forms}

\subsubsection*{Error Function}
\[
\mathrm{erf}(x)=\frac{2}{\sqrt{\pi}}\int_0^x e^{-t^2}\,dt.
\]

\subsubsection*{Complementary Error Function}
\[
\mathrm{erfc}(x)=1-\mathrm{erf}(x)
=\frac{2}{\sqrt{\pi}}\int_x^\infty e^{-t^2}\,dt.
\]

\subsubsection*{$\Phi$-Function (CDF of Standard Normal)}
\[
\Phi(x)=\frac{1}{\sqrt{2\pi}}
\int_{-\infty}^{x}\exp\!\left(-\frac{t^2}{2}\right)dt.
\]

\subsubsection*{$Q$-Function (Gaussian Tail Function)}
\[
Q(x)=\frac{1}{\sqrt{2\pi}}
\int_x^{\infty}\exp\!\left(-\frac{t^2}{2}\right)dt.
\]

\textbf{Note:}
\[
Q(x)=1-\Phi(x).
\]

\subsection*{8.4 Relation Between $\Phi$-Function and $Q$-Function}

\subsubsection*{Evaluating the CDF}

For $X\sim\mathcal N(m,\sigma^2)$,
\[
F_X(x)=\int_{-\infty}^{x}
\frac{1}{\sqrt{2\pi\sigma^2}}
\exp\!\left(-\frac{(y-m)^2}{2\sigma^2}\right)dy.
\]

With substitution $t=\frac{y-m}{\sigma}$,
\[
F_X(x)=\Phi\!\left(\frac{x-m}{\sigma}\right).
\]

\subsubsection*{Evaluating Tail Probabilities}

\[
\Pr(X>x)=\int_x^{\infty}
\frac{1}{\sqrt{2\pi\sigma^2}}
\exp\!\left(-\frac{(y-m)^2}{2\sigma^2}\right)dy
=Q\!\left(\frac{x-m}{\sigma}\right).
\]

\subsubsection*{Key Identities}

\[
Q(x)=1-\Phi(x),
\]
\[
F_X(x)=1-Q\!\left(\frac{x-m}{\sigma}\right).
\]

\textbf{Lecture remarks.}
To evaluate the CDF of a Gaussian RV, evaluate the $Q$-function at the points.
The $Q$-function is more natural for probabilities of the form $\Pr(X>x)$.
$\Phi(x)$ gives area under the left tail; $Q(x)$ gives area under the right tail.

\subsection*{8.5 Gaussian Example}

Given
\[
f_X(x)=\frac{1}{\sqrt{8\pi}}
\exp\!\left(-\frac{(x+3)^2}{8}\right).
\]

Compare with
\[
f_X(x)=\frac{1}{\sqrt{2\pi\sigma^2}}
\exp\!\left(-\frac{(x-m)^2}{2\sigma^2}\right).
\]

Identify:
\[
m=-3,\qquad \sigma^2=4,\qquad \sigma=2.
\]

Gaussian relations:
\[
F_X(x)=\Phi\!\left(\frac{x-m}{\sigma}\right),\qquad
\Pr(X>x)=Q\!\left(\frac{x-m}{\sigma}\right).
\]

\subsubsection*{Probabilities}

\paragraph{1.}
\[
\Pr(X\le 0)=F_X(0)
=\Phi\!\left(\frac{0-m}{\sigma}\right)
=\Phi\!\left(\frac{3}{2}\right)
=1-Q\!\left(\frac{3}{2}\right).
\]

\paragraph{2.}
\[
\Pr(X>4)=Q\!\left(\frac{4-m}{\sigma}\right)
=Q\!\left(\frac{7}{2}\right).
\]

(Other parts are obtained by converting to standard form and expressing answers using
$\Phi(\cdot)$ and $Q(\cdot)$.)

\subsection*{8.6 Types of Continuous Random Variables}

\begin{itemize}
\item Uniform Random Variable (example)
\item Exponential Random Variable (example)
\item Laplace Random Variable (example)
\item Gamma Random Variable
  \begin{itemize}
  \item Graph and special cases
  \item Example
  \end{itemize}
\end{itemize}

Followed by:
\begin{itemize}
\item Homework problem
\item Problem solving
\item In-class activity: Gaussian density estimation
\end{itemize}

\end{document}
